\begin{table}[htbp]
  \centering
  \caption{同一批数据、两个目标量：目标 A（$I_{\mathrm{mwr}}$ vs.\ $I_{\mathrm{rand}}$）与目标 B（$I_{\mathrm{mwr}}$ vs.\ $I_{\mathrm{mix3}}$）}
  \label{tab:targets}
  \footnotesize
  \begin{tabular}{rrrrrrrr}
    \toprule
    \multicolumn{5}{c}{\textbf{目标 A：有好起点值多少}} & \multicolumn{3}{c}{\textbf{目标 B：MWR 相对论文口径的净增量}} \\
    \cmidrule(lr){1-5}\cmidrule(lr){6-8}
    实例 & $\Delta$HV(\%) & 胜出 & $d_z$ & $p$ & $\Delta$HV(\%) & 胜出 & $p$ \\
    \midrule
    Mk01 & +0.047 & 13/30 & +0.06 & $0.685$ & −0.073 & 15/30 & $0.626$ \\
    Mk02 & +0.525 & 22/30 & +0.74 & $5.5\times 10^{-4}{}^{***}$ & +0.010 & 17/30 & $0.556$ \\
    Mk03 & +3.879 & 30/30 & +3.88 & $<10^{-4}{}^{***}$ & +0.023 & 17/30 & $0.328$ \\
    Mk04 & +0.467 & 21/30 & +0.52 & $0.013{}^{*}$ & −0.224 & 12/30 & $0.171$ \\
    Mk05 & +0.238 & 16/30 & +0.24 & $0.184$ & +0.024 & 15/30 & $1.000$ \\
    Mk06 & +4.871 & 30/30 & +5.76 & $<10^{-4}{}^{***}$ & +0.730 & 29/30 & $<10^{-4}{}^{***}$ \\
    Mk07 & +1.593 & 30/30 & +3.25 & $<10^{-4}{}^{***}$ & −0.009 & 16/30 & $1.000$ \\
    Mk08 & +0.944 & 29/30 & +1.55 & $<10^{-4}{}^{***}$ & +0.542 & 28/30 & $<10^{-4}{}^{***}$ \\
    Mk09 & +6.584 & 30/30 & +5.88 & $<10^{-4}{}^{***}$ & +0.012 & 18/30 & $0.685$ \\
    Mk10 & +8.692 & 30/30 & +8.25 & $<10^{-4}{}^{***}$ & +0.690 & 29/30 & $<10^{-4}{}^{***}$ \\
    \bottomrule
  \end{tabular}
  \par\vspace{3pt}\footnotesize\begin{minipage}{\linewidth}$n=\SetSeeds$ seeds，同 seed 配对；$\Delta$HV 为实例边界口径的相对增幅。目标 A 上 10/10 为正（“好起点”这一级确有普遍收益），但\textbf{目标 B 只有 3/10 显著、其余 7 个与 0 不可区分}（配对 Wilcoxon $p\ge0.17$；注意 $\Delta$HV 并非精确的 0，非显著实例中仍为正的最大者是 $+0.024\%$，故“有增益”必须按显著性而非非零来判）：MWR 的净增量是全有或全无，并非“普遍成立但幅度较小”。拿 A 的结论回答 B 的问题即为目标量错位（缺陷 26）。\end{minipage}
\end{table}
